% questions.tex —— 1977 年·北京卷·理科 · 题目内容片段（无答案）
% 由 tools/generate_projects.py 依来源总账自 1977/1977_数学真题_北京理_试题.md 生成，勿手改（改源 MD 后重跑）


\begin{gktip}
\textbf{来源说明}：本题面依据固定源 PDF《1977 北京理科数学试题》第 1 页逐题转录。源文件为 A3 单页扫描，页脚印有“1”；原卷未附答案页。配套答案与解析为本库编者独立推导，\textbf{非官方标准答案}。第 11、12 题按原卷“附加题”保留在第 10 题之后，不计入正式题量。
\end{gktip}

\par\medskip

\begin{question}
\gkti{1}解方程：

\nopagebreak[4]
\[
\sqrt{x-1}=3-x.
\]

\end{question}
\begin{question}
\gkti{2}计算：

\nopagebreak[4]
\[
2^{-\frac12}+\frac{2^0}{\sqrt2}+\frac1{\sqrt2-1}.
\]

\end{question}
\begin{question}
\gkti{3}已知 $\lg2=0.3010$，$\lg3=0.4771$，求 $\lg\sqrt{45}$。

\end{question}
\begin{question}
\gkti{4}证明：

\nopagebreak[4]
\[
(1+\tan\alpha)^2=\frac{1+\sin2\alpha}{\cos^2\alpha}.
\]

\end{question}
\begin{question}
\gkti{5}求过两直线 $x+y-7=0$ 和 $3x-y-1=0$ 的交点且过 $(1,1)$ 点的直线方程。

\end{question}
\begin{question}
\gkti{6}某工厂今年七月份的产值为 $100$ 万元，以后每月产值比上月增加 $20\%$，问今年七月份到十月份总产值是多少？

\end{question}
\begin{question}
\gkti{7}已知二次函数 $y=x^2-6x+5$。

（1） 求出它的图象的顶点坐标和对称轴方程；

（2） 画出它的图象；

（3） 分别求出它的图象和 $x$ 轴、$y$ 轴的交点坐标。

\end{question}
\begin{question}
\gkti{8}一只船以 $20$ 海里/小时的速度向正东航行，起初船在 $A$ 处看见一灯塔 $B$ 在船的北 $45^\circ$ 东方向，一小时后船在 $C$ 处看见这个灯塔在船的北 $15^\circ$ 东方向，求这时船和灯塔的距离 $CB$。

\end{question}
\begin{question}
\gkti{9}一个圆内接三角形 $ABC$，$\angle A$ 的平分线交 $BC$ 于 $D$，交外接圆于 $E$，求证：

\nopagebreak[4]
\[
AD\cdot AE=AC\cdot AB.
\]

\end{question}
\begin{question}
\gkti{10}当 $m$ 取哪些值时，直线 $y=x+m$ 与椭圆

\nopagebreak[4]
\[
\frac{x^2}{16}+\frac{y^2}{9}=1
\]

有一个交点？有两个交点？没有交点？当它们有一个交点时，画出它的图象。

\begin{gktip}
\Needspace{0.25\textheight}
\textbf{附加题}
\end{gktip}

\Needspace{5\baselineskip}
\textbf{附加题 11．}

\nopagebreak[4]
\[
f(x)=
\begin{cases}
x^2\sin\dfrac\pi x,&x\ne0,\\
0,&x=0,
\end{cases}
\]

求函数 $f(x)$ 的导数。

\textbf{附加题 12．}

\begin{itemize}
\item \textbf{（1）} 试用 $\varepsilon-\delta$ 语言叙述“函数 $f(x)$ 在点 $x=x_0$ 处连续”的定义。
\end{itemize}

\begin{itemize}
\item \textbf{（2）} 试证明：若 $f(x)$ 在点 $x=x_0$ 处连续，且 $f(x_0)>0$，则存在一个 $x_0$ 的邻域 $(x_0-\delta,x_0+\delta)$，在这个邻域内，处处有 $f(x)>0$。
\end{itemize}

\end{question}